% !PRINTED 651
\weberpage{651}
12. \gesperrt{Every integral functional $\bar\omega$ is the greatest common
 divisor of two functions $\alpha, \beta$ in $\Omega$.}

To prove this, we take $\alpha = \bar\omega \mu$, divisible by $\bar\omega$,
then $\beta$ divisible by $\bar\omega$ and $\beta : \bar\omega$ relatively prime to $\mu$; then
\[
\bar\omega = \alpha x + \beta y
\]
is the greatest common divisor of $\alpha$ and $\beta$.

\begin{center}
\textbf{§ 183.\qquad Bases and Basis Forms of the Functionals.}
\end{center}

Now let
\weq{1}{\omega_1, \omega_2, \ldots, \omega_n}
be a minimal basis of the field $\Omega$ (§ 176) and
\weq{2}{\alpha_1, \alpha_2, \ldots, \alpha_n}
another basis of $\Omega$. Let the linear substitution by which
the $\alpha$ are connected with the $\omega$ be:
\weq{3}{\alpha_s = \overset{i}{\Sigma} a_{s,i} \omega_i,}
in which the $a_{s,i}$ are integral functions in $Z$, whose determinant
\weq{4}{A = \Sigma \pm a_{1,1} a_{2,2} \ldots a_{n,n}}
does not vanish identically. If $t_1, t_2, \ldots, t_n$ are variables, then
\weq{5}{\lambda = \alpha_1 t_1 + \alpha_2 t_2 + \cdots + \alpha_n t_n}
is the greatest common divisor of $\alpha_1, \alpha_2, \ldots, \alpha_n$.

If one substitutes integral functions in $Z$ for $t_1, t_2, \ldots, t_n$, there arise
 from $\lambda$ integral functions in $\Omega$, all of which are divisible by $\lambda$.

1. \gesperrt{One calls $\lambda$ a basis form and $\alpha_1, \alpha_2, \ldots, \alpha_n$
a basis of the functional $\lambda$, if one obtains all numbers
in $\Omega$ divisible by $\lambda$ by putting
integral functions in $Z$ for $t_i$.}

If the $\alpha_i$ form such a basis, then one can determine integral functions,
$g_{r,s}^{(i)}$ in $Z$, so that
\weq{6}{\alpha_r \omega_s = \overset{i}{\Sigma} g_{r,s}^{(i)} \alpha_i}
holds.

From this it follows:
\weq{7}{\lambda \omega_s = \overset{i}{\Sigma} \alpha_i t_{i,s},}
if
\weq{8}{t_{i,s} = \overset{r}{\Sigma} g_{r,s}^{(i)} t_r}
%CONT>>
